\documentclass[12pt]{article}
\usepackage[margin=2.5cm]{geometry}
\usepackage[spanish]{babel}
\usepackage{graphicx}
\usepackage{parskip}
\usepackage{array}
\renewcommand{\baselinestretch}{1.25}
\pagestyle{empty}

\begin{document}

% === PÁGINA 1: RESOLUCIÓN PRINCIPAL ===
\begin{minipage}{0.3\textwidth}
  \vspace{-0.6cm}
  % RUTA CORREGIDA: Le decimos que suba un nivel (../) y entre a la carpeta de recursos.
  \includegraphics[width=2.5cm]{../05_Templates_y_Recursos/logo.png}
\end{minipage}
\hfill
\begin{minipage}{0.65\textwidth}
  \begin{flushright}
    Buenos Aires, 24 de diciembre de 2025
  \end{flushright}
\end{minipage}

\vspace{1cm}

\begin{center}
  \textbf{\large Resolución r24eVIIs25 – Presupuesto mensual de diciembre}
\end{center}

\vspace{0.8cm}

\textbf{VISTO:}


\vspace{0.4cm}

\textbf{CONSIDERANDO:}
\begin{itemize}
  \item Que la mensualidad correspondiente al mes de noviembre fue de \textbf{\$144.999 ARS}.
  \item Que dicho monto fue utilizado en los siguientes conceptos:
    \begin{itemize}
      
      \item Transporte: \textbf{\$14.000}
      
      \item Comida: \textbf{\$27.000}
      
      \item Regalo (cumpleaños): \textbf{\$26.000}
      
      \item Celular (cuota 1/12): \textbf{\$42.416}
      
      \item Peluquería: \textbf{\$10.000}
      
      \item Inversiones: \textbf{\$16.000}
      
    \end{itemize}
  
  \item Que para el mes de julio se proyecta un presupuesto total de \textbackslash{}textbf\{\textbackslash{}\$131.416 ARS\}, eliminando el rubro “Ocio y otros” como medida de ajuste.
  
  \item Que el solicitante mantiene un compromiso mensual de pago por el celular POCO X6 Pro 5G, abonando la cuota \textbackslash{}textbf\{2/12\} de \textbackslash{}textbf\{\textbackslash{}\$42.416 ARS\}.
  
  \item Que el solicitante ha recibido una transferencia parcial de \textbackslash{}textbf\{\textbackslash{}\$50.000 ARS\} correspondiente a parte del presupuesto.
  
  \item Que se debe aplicar una penalización de \textbackslash{}textbf\{\textbackslash{}\$2.500 ARS\} por llegada tarde registrada durante junio.
  
  \item \$1.200 aguante eñ \%es
  
\end{itemize}

\vspace{0.4cm}

\textbf{RESUELVO:}

\textbf{ARTÍCULO 1°.-}
Aprobar el presupuesto mensual correspondiente al mes de julio de 2025 por un total de \textbackslash{}textbf\{\textbackslash{}\$131.416 ARS\}, conforme al detalle del Anexo I.
\vspace{0.5cm}

\textbf{ARTÍCULO 2°.-}
Registrar que: \textbackslash{}begin\{itemize\} \textbackslash{}item Se ha transferido previamente la suma de \textbackslash{}textbf\{\textbackslash{}\$50.000 ARS\}. \textbackslash{}item La cuota \textbackslash{}textbf\{2/12\} del celular (\textbackslash{}textbf\{\textbackslash{}\$42.416 ARS\}) será abonada directamente por el solicitante y no será transferida. \textbackslash{}item La penalización por llegada tarde (\textbackslash{}textbf\{\textbackslash{}\$2.500 ARS\}) será descontada del monto final. \textbackslash{}end\{itemize\} Por lo tanto, el monto neto a transferir asciende a \textbackslash{}textbf\{\textbackslash{}\$36.500 ARS\}.
\vspace{0.5cm}

\textbf{ARTÍCULO 3°.-}
El solicitante continuará registrando y rindiendo cuentas de sus gastos mensuales, conforme al esquema de control financiero vigente.
\vspace{0.5cm}


\vspace{1.4cm}

\begin{flushright}
  % RUTA CORREGIDA: También para la firma.
  \includegraphics[width=3.2cm]{../05_Templates_y_Recursos/firma.png}
\end{flushright}

% === PÁGINA 2: ANEXO I (SI EXISTE) ===


\end{document}